% !TeX root = ../papers/nc-p3.tex
\section{有理正規三次上の基点を持つ階数二関係}
\label{sec:twisted-cubic-line-modules}
各辺の両射影のファイバーが空か一点である標準有限Frobenius窓で、
点列の中に全位置で完全な三次系により埋め込まれる同じ $C=\mathbb P^1$ があるとする。
曲線上の点の移動に沿って座標を取り、$V_i=S^3$ と書く（$S^m=H^0(C,\mathcal O(m))$）。
$H_i(p)$ は点 $p$ の消滅超平面である。曲線外の有限構造には条件を付けない。
\begin{theorem}\label{thm:twisted-cubic-line-resolution}
$r\in R_i$ が階数2で、その左支持 $U\subset V_i$ の二つの三次式が共通零点を持つとする。
右支持を $W$ とする。右加群 $L_i(U)=e_i A/U A$ の長さ $n$ の次元は全て $n+1$ であり、
$U\otimes e_{i+1}A\to e_iA$ の関係は $r$ 一個で生成される。
左加群についても両端を反転した結論を得る。
さらにこの窓が標準Koszul左右AS正則であることと、全始点・長さで
$\dim A_{i,i+n}=\binom{n+3}{3}$ となることは同値である。
\end{theorem}
\begin{proof}[点加群による線形表示]
$r$ の両支持の鉛筆は基因子を除くと同じ次数の射を与える。
実際、二つの独立な左係数 $u,v$ と右係数 $s,t$ で $r=u\otimes t-v\otimes s$ と書くと、
対角上で $ut=vs$ である。従って $u=dF,v=dG,s=gF,t=gG$ と書け、
$F,G$ は互いに素で $\deg d=\deg g$ である。特に左右の基点の有無は一致する。
非零の階数1関係は積が整曲線上で消えないため存在しない。

共通零点 $p$ を一つ取り、$q\in\mathbb P(U^\perp)$ を $p$ と異なる点とする。
$H_i(p)=U\oplus kv$ として $q(v)=1$ とする。
右点加群 $M_i(p)=e_iA/H_i(p)A$ のこの超平面の関係は
$K=R_i\cap(H_i(p)\otimes V_{i+1})$ の三個で生成される。
これは命題\ref{prop:commuting-three-point-graph}の点列での階数 $1,3,3,1$ の計算である。
$p$ の前点で四テンソルを縮約した三テンソルを $h$ とすれば、その末尾縮約が $K$ を張る。
$h$ の先頭を $q$ で縮約した $r'$ は零でない。
実際、前点から出る辺の六関係上の対評価が零なら、点ファイバーの一意性で $q=p$ となる。
相補二位置との完全対によりその対評価はこの縮約で読める。
$K$ での先頭 $q$ 評価は $r$ を消すので、その像 $U'$ は次元2以下であり、
$r'\in R_{i+1}$ の左支持である。階数1の関係がないので次元は2、核は $kr$ となる。
$h$ の最終支持は四移動で同一視した前点の消滅超平面に入る（巡回した隣接グラフでの消滅）。
従って $r'$ の右支持も曲線上に基点を持ち、先の分解により $U'$ にも基点がある。

$va\in UA$ は $H_i(p)$ の関係を意味し、その $v$ 成分全体は上の三生成と $q$ 評価により
ちょうど $U'A$ である。従って次数を各始点から測ると
\[
 0\longrightarrow L_{i+1}(U')(-1)\xrightarrow{\ v\ }
 L_i(U)\longrightarrow M_i(p)\longrightarrow0 .
\]
全ての後続加群でこの操作を反復し、長さ0から帰納すれば主張の次元を得る。

$K$ の基底を $r,s_1,s_2$ とし、後二者の $v$ 成分を $u_1',u_2'$ とする。
$h=r\otimes c_0+s_1\otimes c_1+s_2\otimes c_2$ と書けば
$r'=u_1'\otimes c_1+u_2'\otimes c_2$ であり、
$h\in H_i(p)\otimes R_{i+1}$ により $r c_0+s_1 c_1+s_2 c_2=0$ が $H_i(p)\otimes e_{i+1}A$ で成り立つ。
$U$ の任意の関係を $r a_0+s_1 a_1+s_2 a_2$ と表示すると
$u_1'a_1+u_2'a_2=0$ である。一つ短い長さでの帰納仮定により
$(a_1,a_2)=(c_1 b,c_2 b)$ だから、元の関係は $r(a_0-c_0 b)$ である。
開始次数2では $q$ 評価の核 $kr$ そのものとなる。左側でも逆向きの点列で同じ計算ができる。
\end{proof}

\begin{proof}[Hilbert関数からの十分性]
どの位置でも任意の $p\in C$ を含む二次有効因子 $D$ で
$U=H^0(\mathcal O(3)(-D))\subset H_i(p)$ を左支持とする階数2関係が存在する。
実際、$R_i=[z,t]T_i$ と二次の六平面 $T_i\subset S^2\otimes S^2$ を書き、
$D$ の式を $d$ とすると、求める条件はある $g\ne0$ に対する $d\otimes g\in T_i$ である。
$T_i^\perp$ の三行を $d$ で評価した三次行列式は、
$p$ を根に持つ $d$ の射影直線で必ず零点を持つ。左加群用の逆向きでも同じである。

所定のHilbert関数を仮定する。上で得た三項線形表示のEuler式から、
左端の $e_{i+2}A\to U\otimes e_{i+1}A$ も単射となる。
左表示にも適用できる。右表示の双対は、階数2の係数対の基底を合わせると
同じ $r$ の右支持 $W$ から作る左表示の末尾の商を除いた複体である。
従って各右 $L_i(U)$ から自由加群への次数0,1のExtは消える（各表現可能加群を標的とする）。
任意の点加群は先に構成したこの種の $L$ の商なので自由加群へのHomは零であり、その尾部も点加群である。
上の短完全列を正次数に制限し、$L_{i+1}(U')$ と点加群の尾部の両方でHomが消えることから
$L_i(U)_{\ge1}$ のHomも消える。

右単純の第一Extの消滅を見るため、$e_iA_+$ から自由加群への任意の写像を $UA$ に制限する。
$L_i(U)$ の第一Ext消滅でこれは $e_iA$ 上に延びる。差は $L_i(U)_{\ge1}$ を通るので零である。
単純からのHomも、$L_i(U)$ の商であるため零である。左側でも同じで、
候補Koszul複体の生成・二次表示の端とそのFrobenius双対の端が完全になり、中央はEuler式による。
逆向きは標準Koszul解消のEuler式である。
\end{proof}

\begin{corollary}\label{cor:twisted-cubic-compatible-pencils}
全添字で二次元鉛筆 $U_i\subset V_i$ があり、
$R_i\cap(U_i\otimes U_{i+1})$ が基点を持つ階数2関係 $r_i$ を含めば、この窓は標準正則である。
特に四移動の任意の指定した有限べきで閉じるこのような列は有限方程式による十分条件を与える。
\end{corollary}
\begin{proof}
二文字・一非退化関係の添字付き代数 $P$ は、基底を両方向に順に合わせれば二変数多項式環の反復である。
定理の関係生成を $P$ 単純のKoszul複体に適用すると $\operatorname{Tor}_1^P(k_i,A)=0$。
各終点で有界下加群の極小表示を取ると $A$ は左 $P$ 自由である。自由生成元の個数は
$e_iA/U_iA$ の次元から長さ $n$ で $n+1$ なので、Hilbert級数は $(1-t)^{-4}$ となる。
\end{proof}

\begin{proposition}\label{prop:twisted-cubic-secant-cubic}
$T_i^\perp$ の $d$ 評価の行列式が恒等零でなければ、その平面三次式は必ず直線因子を持つ。
$g$ 側の行列式についても同じである。
\end{proposition}
\begin{proof}
$d$ の根を順序付きで $p,q$ とする。$q\ne p$ では $q$ の曲線上の像を
上記の $q$ 評価に使い、$q=p$ では接線上の微分評価を使う。
このとき $H_i(p)$ 内の評価核はいずれも $d S^1$ である。従って求める三次式の消滅は、
三テンソルの先頭をこの評価で縮約した行列の階数が2以下であることに等しい。
四テンソルを四辺の角括弧と多重一次式 $f$ の積に書く。
その順序を $p,q,z,t$ とすると、$p,q$ での縮約から $[p,q]$ と両端の
一次式（残った二因子に作用する単射）を除いたものは $[z,t] f(p,q,z,t)$ である。
継ぎ目の四移動は外側の一次式か $p$ の座標同一視に繰り入れる。
$q=p$ では $q$ 側の微分が消えた括弧にだけ作用し、同じ行列を与える。

$f(p,q,z,t)=a+bz+ct+d'tz$ と書くアフィン座標では、残った三行三列の行列式は
$(c-b)(ad'-bc)$ である。ここで $c-b$ は $p,q$ について二次数 $(1,1)$ の切断である。
これが恒等零なら元の三次式も恒等零である。そうでなければ、この切断の零因子の
各既約成分の対称平方平面での像は直線または二次曲線である。
実際、順序を忘れる有限写像による超平面束の引戻しは $\mathcal O(1,1)$ で、
その各成分上の次数は1または2である。元の三次式はこの像で消える。
二次曲線を含む場合も三次式の残余は直線である。両側の交換で他方の結論を得る。
\end{proof}
